\documentclass[a4paper,12pt]{article}
\usepackage[utf8]{inputenc}
\usepackage{xeCJK}
\usepackage{fontspec}
\usepackage{geometry}
\usepackage{titlesec}
\usepackage{indentfirst}
\usepackage{enumitem}
\usepackage{fancyhdr}
\usepackage{graphicx}
\usepackage{xcolor}
\usepackage{tocloft}
\usepackage{setspace}
\usepackage{hyperref}
\usepackage{tikz}
\usepackage{booktabs}
\usepackage{array}
\usepackage{tabularx}
\usepackage{float}
\usepackage{longtable}
\usepackage{caption}
\usepackage{amsmath}
\usepackage{amssymb}
\usepackage{listings}
\usepackage[most]{tcolorbox}
\usetikzlibrary{arrows.meta,positioning,shapes.geometric}

\geometry{left=2.5cm, right=2.5cm, top=3cm, bottom=2.5cm, headheight=24.01996pt}
\addtolength{\topmargin}{-0.01996pt}

\setmainfont{Times New Roman}
\IfFontExistsTF{Hiragino Sans GB}{
  \setCJKmainfont{Hiragino Sans GB}
  \setCJKmonofont{Hiragino Sans GB}
}{
  \IfFontExistsTF{Microsoft YaHei}{
    \setCJKmainfont{Microsoft YaHei}
    \setCJKmonofont{Microsoft YaHei}
  }{
    \setCJKmainfont{SimSun}
    \setCJKmonofont{SimSun}
  }
}

\definecolor{deepblue}{RGB}{0,51,153}
\definecolor{softblue}{RGB}{51,102,153}
\definecolor{codebg}{RGB}{247,249,252}
\definecolor{codeframe}{RGB}{163,181,214}
\definecolor{codecomment}{RGB}{76,114,54}
\definecolor{codestring}{RGB}{180,102,0}

\lstdefinestyle{pyreport}{
  language=Python,
  basicstyle=\ttfamily\small,
  keywordstyle=\color{deepblue}\bfseries,
  commentstyle=\color{codecomment}\itshape,
  stringstyle=\color{codestring},
  numberstyle=\scriptsize\color{black!45},
  numbers=left,
  numbersep=8pt,
  breaklines=true,
  breakatwhitespace=false,
  showstringspaces=false,
  keepspaces=true,
  tabsize=4,
  columns=fullflexible,
  frame=none
}

\newtcblisting{codeblock}[1][]{
  enhanced,
  breakable,
  listing only,
  listing engine=listings,
  colback=codebg,
  colframe=codeframe,
  boxrule=0.6pt,
  arc=1.5mm,
  left=1.2mm,
  right=1.2mm,
  top=1.0mm,
  bottom=1.0mm,
  listing options={style=pyreport},
  #1
}

\renewcommand{\cftsecfont}{\bfseries\color{softblue}}
\renewcommand{\cftsubsecfont}{\color{black}}
\renewcommand{\cftsubsubsecfont}{\color{black}}
\renewcommand{\cftsecpagefont}{\color{softblue}}
\renewcommand{\cftsubsecpagefont}{\color{black}}
\renewcommand{\cftsubsubsecpagefont}{\color{black}}
\setlength{\cftbeforesecskip}{10pt}
\setlength{\cftbeforesubsecskip}{5pt}
\setlength{\cftsecindent}{0em}
\setlength{\cftsubsecindent}{2em}
\renewcommand{\cftsecdotsep}{\cftdotsep}
\renewcommand{\cftsubsecdotsep}{\cftdotsep}
\renewcommand{\cftdot}{.}
\renewcommand{\cfttoctitlefont}{\huge\bfseries\color{softblue}}
\renewcommand{\contentsname}{目录}

\titleformat{\section}
  {\color{deepblue}\large\bfseries}
  {\color{deepblue}\thesection}
  {0.5em}
  {}
  [{\vspace{-6pt}\color{deepblue}\rule{\textwidth}{0.8pt}\vspace{3pt}}]

\titleformat{\subsection}
  {\normalsize\bfseries}
  {\thesubsection}
  {0.5em}
  {}

\titleformat{\subsubsection}
  {\normalsize\bfseries}
  {\thesubsubsection}
  {0.5em}
  {}

\pagestyle{fancy}
\fancyhf{}
\fancyhead[L]{\small\color{deepblue} 计算机视觉工业视频缺陷检测实验报告}
\fancyhead[R]{\small\color{deepblue} Side Micro ROI Detector}
\fancyfoot[R]{\small\thepage}
\renewcommand{\headrulewidth}{0.4pt}
\renewcommand{\footrulewidth}{0pt}
\renewcommand{\headrule}{\hbox to\headwidth{%
  \color{deepblue}\leaders\hrule height \headrulewidth\hfill}}

\onehalfspacing
\setlength{\parindent}{2em}
\setlength{\parskip}{0pt}
\emergencystretch=2em
\graphicspath{{../}}

\newcolumntype{Y}{>{\raggedright\arraybackslash}X}
\newcolumntype{L}[1]{>{\raggedright\arraybackslash}p{#1}}
\newcommand{\code}[1]{\texttt{#1}}

\hypersetup{
    colorlinks=true,
    linkcolor=softblue,
    filecolor=magenta,
    urlcolor=blue,
    citecolor=blue,
    pdftitle={面向工业视频的侧边接口套缺失检测实验报告},
    pdfauthor={Project Report Draft},
    pdfsubject={工业视频缺陷检测、固定微小 ROI、Patch 分类与视频事件统计},
    pdfkeywords={工业缺陷检测, 视频理解, 小样本 ROI 校准, 微小 ROI, ExtraTrees, 事件聚合}
}

\begin{document}

\begin{titlepage}
    \noindent\makebox[\linewidth]{\color{deepblue}\rule{\paperwidth}{0.2cm}}
    \vspace*{1cm}

    \begin{center}
        {\LARGE\bfseries 计算机视觉项目实验报告\par}
        \vspace{1.2cm}
        {\Huge\bfseries 面向工业视频的侧边接口套缺失检测\par}
        \vspace{0.6cm}
        {\Large\bfseries 从自动大 ROI 失败到固定微小 ROI Patch 分类与视频事件统计\par}
        \vspace{0.3cm}
        {\large 基于微小 ROI 校准、上下水平检测带、规则增强与可视化回放的工程迭代报告\par}
    \end{center}

    \vspace{2.2cm}

    \begin{flushright}
        {\large 姓名：严广粟\par}
        \vspace{0.3cm}
        {\large 班级：23计科02\par}
        \vspace{0.3cm}
        {\large 指导老师：金长龙\par}
        \vspace{0.3cm}
        {\large GitHub：\url{https://github.com/ximiximixi/work3}\par}
        \vspace{0.3cm}
        {\large 报告日期：2026年5月16日\par}
    \end{flushright}

    \vspace{2.2cm}
    \noindent\makebox[\linewidth]{\color{deepblue}\rule{\paperwidth}{0.2cm}}
\end{titlepage}

\pagenumbering{roman}

\begin{abstract}
\color{black}
\setlength{\parindent}{2em}
\indent 本项目面向工业检测视频中的侧边接口装配缺陷识别任务，目标是在视频帧中判断目标样本侧边接口区域应出现的米白色套/盖是否缺失。该问题并不是普通的图像分类任务：缺陷证据发生在右半区域内极小的侧边接口 patch 中，且正常米白色实体、透明接口、金属反光和背景结构之间存在明显视觉混淆。项目早期曾尝试自动定位较大的目标 ROI 或整根样本区域，并结合纹理异常检测、颜色证据和端盖定位进行判别，但大 ROI 会混入背景、透明件、金属反光和无关结构，使正常/缺陷判断不稳定。该失败结果推动项目从“让算法自动猜大区域”转向“重新定义真正需要看的小检测区域”。

\indent 在用户进一步指出缺陷只发生于侧边很小接口区域后，项目构建了一个小规模 ROI 校准样本集。ROI 校准工具仅展示原视频右半边裁剪图，并将 ROI 坐标加回 \code{offset\_x=960} 保存为原始视频坐标。最终校准集包含 $22$ 帧、$115$ 个 \code{sample} 小框，其中 \code{PRESENT}=99、\code{MISSING}=16；上侧检测带 $58$ 个、下侧检测带 $57$ 个。ROI 样本分析表明，\code{sample} 框并非整根样本，而是实际需要检测的 side micro ROI。因此最终算法采用固定上下两个水平检测带，上侧约为 $y=350$--$495$，下侧约为 $y=835$--$960$，并在蓝色端盖/结构线索约束下生成少量候选 patch。

\indent 最终实现位于 \code{side\_micro\_roi\_detector.py}。模型以 ROI 校准小框和随机背景 patch 训练三分类 \code{PRESENT/MISSING/BACKGROUND} ExtraTrees 分类器，并使用颜色、亮度、HSV/Lab 统计、空间池化、米白色区域比例、边缘密度和纹理直方图等 $127$ 维特征。为修复透明接口被误判为正常的问题，检测器进一步加入 \code{white\_ratio}、\code{sleeve\_blob} 与 \code{edge\_density} 规则，强调“米白色实体套/盖证据”而不是透明接口本身。最终全视频以 $0.5$s 间隔扫描，输出 $12{,}762$ 条 patch 观测、$92$ 个持续事件；最后一分钟 showcase 以 $8$ fps 渲染 $481$ 帧，输出 $30$ 个事件，并生成包含风险仪表、CLEAR/ALERT 状态、风险曲线、上下通道风险条、证据 crop 和热力图的高级可视化视频。由于当前项目尚无独立逐帧或事件级 ground truth，本报告将这些结果解释为工程输出统计与定性验证基础，而不将其夸大为严格准确率结论。
\end{abstract}

\vspace{1cm}
\begin{center}
    {\large\bfseries\color{black} 关键词：\color{black}工业缺陷检测；视频理解；微小 ROI；ROI 校准；Patch 分类；事件聚合；可视化回放}
\end{center}

\clearpage
\tableofcontents
\clearpage

\pagenumbering{arabic}
\color{black}

\section{引言}
\subsection{项目背景}
工业视频检测项目往往不同于标准公开数据集上的图像分类。公开数据集通常假设目标对象已经裁剪、居中且类别边界明确；而实际产线视频中的缺陷判断首先依赖正确的观测区域定义。本项目关注的对象是工业视频中目标样本侧边接口位置的米白色套/盖是否缺失。正常状态下，该位置应具有明显的米白色实体套/盖证据；缺陷状态下，该实体可能消失，视觉上表现为透明接口、空洞感增强、内部边缘增多、蓝色或背景结构暴露，以及局部纹理和反光变化。

从工程过程看，本项目最重要的经验并不是“训练了一个分类器”，而是通过失败实验和任务反馈逐步发现：真正有效的信息只存在于上、下两个侧边水平窄带中的小 patch，而不是整根样本或较大的自动 ROI。若检测区域过大，算法会同时看到整根滤芯、蓝色端盖、透明件、金属反光、背景干扰和遮挡结构，缺陷证据反而被大量无关视觉信息稀释。

\subsection{报告范围}
本报告基于当前项目文件夹中已有文件整理，包括 \code{paper/} 模板、顶层 Python 脚本、ROI 校准前端、ROI 样本文件、训练与检测输出、可视化图片、结果视频、早期实验目录和旧版本输出。本报告的目标是还原完整实验路线，而不是只描述最终成功版本。因此，报告同时记录以下内容：
\begin{enumerate}[label=\arabic*.]
    \item 早期自动 ROI / 大检测框 / PatchCore 或颜色基线尝试及其局限；
    \item ROI 小样本校准集和校准工具的设计；
    \item 从 ROI 样本中重新理解 side micro ROI 的过程；
    \item 固定上下两个水平检测带 + 小 patch 分类器的最终方法；
    \item 规则增强、事件过滤、批量预测和重叠框筛选修复；
    \item 低帧率检查视频和最后一分钟高级 showcase 视频；
    \item 当前仍缺少独立测试集和事件级真值所带来的解释边界。
\end{enumerate}
\newpage
\section{任务定义}
\subsection{输入与输出}
项目输入为工业检测视频 \code{20241101\_161258.mp4} 及其采样帧。根据 ROI 校准任务文件，该视频帧分辨率为 $1920\times1080$，校准工具仅展示右半边裁剪区域，即原始坐标 $x\geq960$ 的区域。最终检测器进一步将关注范围缩小到两个固定水平带：上侧侧边接口检测带 $y=350$--$495$，下侧侧边接口检测带 $y=835$--$960$。候选 patch 的默认尺寸为 $84\times68$，主要围绕蓝色端盖或结构锚点附近生成。

系统输出包括四个层次：
\begin{enumerate}[label=\arabic*.]
    \item patch 级状态：\code{PRESENT}、\code{MISSING}、\code{BACKGROUND} 或中间的 \code{UNKNOWN}；
    \item 帧级观测：每个采样时刻上、下通道候选 patch 的状态、概率、规则证据和坐标；
    \item 事件级统计：持续缺失事件的开始时间、结束时间、持续观测数、峰值缺失概率和触发通道；
    \item 可视化结果：检测时间线、$2$ fps 检查视频、最后 $60$s 的 $8$ fps showcase 视频及关键帧截图。
\end{enumerate}

\subsection{缺陷定义}
本项目中的缺陷定义为：在目标样本侧边接口区域本应观察到米白色套/盖实体的位置，没有观察到足够明确的米白色实体证据。正常样本不只是“接口可见”，而是应存在连续、低饱和、高亮度、连通性较好的米白色套/盖区域。缺陷样本可能仍然出现透明接口或局部高亮反光，但这些透明结构不能被当作正常证据。

\subsection{任务难点}
该任务具有以下困难：
\begin{enumerate}[label=\arabic*.]
    \item 缺陷区域非常小，常见 patch 仅约 $80\times70$ 像素；
    \item 透明接口会产生类似正常部件的轮廓，容易被误判为 \code{PRESENT}；
    \item 金属反光和背景结构在大 ROI 中占比高，会显著干扰纹理和颜色统计；
    \item 缺陷样本稀缺，ROI 校准集中 \code{MISSING} 仅 $16$ 个原始 patch；
    \item 目标在视频中存在轻微位置变化，因此不能简单地对全宽密集滑窗报警；
    \item 若没有事件级过滤，单帧反光、遮挡或定位抖动会造成孤立误报。
\end{enumerate}

\section{项目资产与模板理解}
\subsection{顶层目录概览}
项目顶层包含两个原始视频、多个阶段性输出目录、最终 side micro ROI 检测输出目录、ROI 校准工具、DINO/颜色基线尝试、PatchCore 参考仓库和报告模板。表~\ref{tab:top-level-index} 给出核心目录功能概览，完整资产索引见附录 A。

\begin{table}[H]
\centering
\caption{项目顶层目录功能概览}
\label{tab:top-level-index}
\small
\begin{tabularx}{\textwidth}{L{0.30\textwidth}Y}
\toprule
路径 & 功能说明 \\
\midrule
\path{cap_annotation_tool/} & ROI 校准前端、任务帧、校准 JSON/CSV、side micro ROI crop 与校准可视化。 \\
\path{side_micro_outputs/} & 最终固定上下水平带 + 小 patch 分类器输出，包括模型、训练 manifest、观测、事件、时间线、检查视频和 showcase 视频。 \\
\path{cap_defect_outputs/} & 自动端盖/较大 ROI 检测尝试，包含蓝色端盖定位、伪正常训练、异常事件与调试图。 \\
\path{defect_outputs/} & 更早期的粗 ROI / PatchCore 风格纹理异常检测输出，包含事件分析和高级诊断图。 \\
\path{dino/} & DINO/颜色基线框架与输出，用于探索通用异常检测与演示视频生成；当前未形成有标签精度结论。 \\
\path{patchcore-inspection-main/} & PatchCore 参考实现和预训练/缓存模型资产，作为早期异常检测思路背景。 \\
\path{paper/} & LaTeX 报告模板目录，当前已有 \path{hpc_ch2.tex}。 \\
\bottomrule
\end{tabularx}
\end{table}




\newpage
\section{技术路线演化}
\subsection{从粗 ROI 检测到微小 ROI 分类}
图~\ref{fig:pipeline-evolution} 总结了本项目的技术路线。该路线并非一次成型，而是在多次失败、需求反馈和可视化检查中逐步收敛。

\begin{figure}[H]
\centering
\resizebox{\textwidth}{!}{%
\begin{tikzpicture}[
    node distance=0.8cm,
    block/.style={rectangle, rounded corners=1.5mm, draw=deepblue, very thick, align=center, fill=blue!4, minimum width=3.4cm, minimum height=1.0cm},
    bad/.style={rectangle, rounded corners=1.5mm, draw=red!65!black, very thick, align=center, fill=red!5, minimum width=3.4cm, minimum height=1.0cm},
    good/.style={rectangle, rounded corners=1.5mm, draw=green!50!black, very thick, align=center, fill=green!5, minimum width=3.4cm, minimum height=1.0cm},
    arrow/.style={-{Latex[length=3mm]}, thick, color=deepblue}
]
\node[bad] (coarse) {自动大 ROI / 粗检测框\\混入背景和反光};
\node[block, right=of coarse] (anno) {ROI 小样本校准集\\22 帧 / 115 patch};
\node[block, right=of anno] (reinterpret) {重新解释 sample 框\\side micro ROI};
\node[good, right=of reinterpret] (micro) {固定上下水平带\\小 patch 分类器};
\node[good, below=of micro] (rule) {米白色实体证据规则\\透明接口不算正常};
\node[good, left=of rule] (event) {帧级观测\\事件级聚合};
\node[good, left=of event] (vis) {2fps 检查视频\\8fps showcase};
\draw[arrow] (coarse) -- (anno);
\draw[arrow] (anno) -- (reinterpret);
\draw[arrow] (reinterpret) -- (micro);
\draw[arrow] (micro) -- (rule);
\draw[arrow] (rule) -- (event);
\draw[arrow] (event) -- (vis);
\end{tikzpicture}
}
\caption{从自动大 ROI 检测失败到固定 side micro ROI 分类与事件可视化的技术路线演化}
\label{fig:pipeline-evolution}
\end{figure}

\subsection{第一阶段：自动 ROI 与大检测框尝试}
项目最初尝试对整根样本、主目标区域或较大的右半区域进行自动检测。对应代码和输出包括 \path{defect_detection_pipeline.py}、\path{defect_outputs/}、\path{cap_defect_detection_pipeline.py}、\path{cap_defect_outputs/} 以及 \path{dino/} 中的颜色基线/AnomalyDINO 框架。早期方法尝试使用圆形/端盖定位、颜色阈值、PatchCore 风格的 patch memory 和异常分数时间线来发现装配异常。

这些实验为项目提供了有价值的失败证据：大检测框能够捕获许多与目标相关的结构，但也不可避免地把透明件、蓝色端盖、金属反光、整根滤芯、背景和遮挡纳入判别区域。对于本任务而言，缺陷证据并不分布在整个样本上，而集中在侧边接口处非常小的米白色套/盖区域。大 ROI 的视觉统计被无关区域主导，使缺失判断不可靠。

\subsection{第二阶段：ROI 小样本校准集构建}
为避免算法继续猜测 ROI，项目构建了一个小规模 ROI 校准样本集。脚本 \path{prepare_cap_annotation_set.py} 从 \code{20241101\_161258.mp4} 中抽取 $22$ 帧右半边裁剪图，覆盖正常参考帧、中段疑似样本和后段稀缺异常样本。特别地，抽帧包含 $1072.0$s 到 $1084.0$s 附近的密集时间点，用于覆盖后段罕见异常。

ROI 校准前端位于 \path{cap_annotation_tool/}，包括 \path{server.py}、\path{index.html}、\path{app.js} 和 \path{style.css}。前端只显示右半边裁剪图，保存时自动加回 \code{offset\_x=960}，使最终坐标位于原始 $1920\times1080$ 视频坐标系中。ROI 样本 schema 支持帧级 \code{NORMAL/DEFECT/UNKNOWN}，部位级 \code{top/bottom/side\_upper/side\_lower} 缺失勾选，以及 \code{sample/fixed\_region/top/bottom/side\_upper/side\_lower/ignore} 等框类型。

\subsection{第三阶段：重新理解 ROI 样本}
读取 \path{annotations.json} 和 \path{side_roi_training_summary.json} 后发现，\code{sample} 框并不是整根样本，而是侧边很小的接口检测区域。该发现是项目的关键转折：算法不应继续检测整根滤芯或较大端盖区域，而应只看上、下两个水平窄带中的小 patch。

统计文件 \path{cap_annotation_tool/side_roi_training_summary.json} 给出的结论非常明确：ROI 小框共 $115$ 个，其中 \code{PRESENT}=99、\code{MISSING}=16；上侧 $58$ 个、下侧 $57$ 个。推荐固定检测带为上侧 $y=350$--$495$、下侧 $y=835$--$960$，crop 尺寸建议为 $96\times76$。最终检测代码采用接近该尺度的 $84\times68$ 窗口。

\subsection{第四阶段：固定水平带 + 小 Patch 分类器}
最终算法新建在 \path{side_micro_roi_detector.py} 中，不覆盖旧的大 ROI 流水线。它直接读取 \path{cap_annotation_tool/tasks.json} 和 \path{cap_annotation_tool/annotations.json}，将 \code{sample} 框解释为 side micro ROI，并构建 \code{PRESENT/MISSING/BACKGROUND} 三分类训练集。候选生成不采用全宽密集滑窗，而是先通过旧端盖检测模块中的蓝色锚点定位目标结构，再围绕锚点中心施加几个固定 $x$ 偏移生成少量候选；若锚点缺失，则退化为较稀疏的右半区域候选。

\subsection{第五阶段：规则增强与误检修复}
训练初版小 patch 分类器后，仍出现两个典型问题。第一，透明接口可能被模型当作正常结构；第二，重叠候选框排序时，旁边的高置信 \code{PRESENT} 框可能压掉真实 \code{MISSING} 框。为此，最终代码增加 \code{sleeve\_evidence} 和 \code{refine\_state\_with\_sleeve\_rule}，把判断重点放在米白色实体证据上：低饱和高亮区域比例 \code{white\_ratio}、最大连通白色实体 \code{sleeve\_blob}、边缘密度 \code{edge\_density}。同时，NMS 替换逻辑允许强 \code{MISSING} 候选替换弱 \code{PRESENT} 候选，也允许确凿实体套/盖证据恢复误判缺失。

\subsection{第六阶段：视频可视化展示}
最终输出不只包含 CSV/JSON，还生成了两个视频资产。第一个是 \path{side_micro_outputs/visuals/02_micro_review_2fps.mp4}，用于全视频低帧率检查；第二个是 \path{side_micro_outputs/visuals/03_last60s_8fps_showcase.mp4}，覆盖最后约 $60$s，使用 $8$ fps 渲染，并在左侧无用区域加入高级动态面板：实时风险仪表、\code{CLEAR/ALERT/WATCH} 状态、最大缺失分数曲线、上下通道风险条、证据 crop、热力图和最近 $20$s 风险条带。该视频是项目展示和复核的重要成果，但由于没有逐帧真值，它应被解释为定性展示，而非严格精度评估。
\newpage
\section{\code{20241031\_222226.mp4} 视频处理实验}
\subsection{实验定位}
\path{20241031_222226.mp4} 是本项目中最早形成完整输出闭环的视频实验。它并不是最终 side micro ROI 检测器的输入视频，但它完整体现了项目第一阶段的技术路线：先在较大的目标区域内定位圆形/端面目标，再用 PatchCore 风格的纹理记忆库和颜色/纹理惩罚项生成异常分数，最后通过事件过滤和回放视频进行视频检查。

该实验的价值在于，它证明了“粗 ROI 异常检测”能够在某些片段上形成连续事件，但也暴露了本项目后续必须转向微小侧边 ROI 的原因：大 ROI 能捕捉整体异常趋势，却无法严格回答“侧边接口米白色套/盖是否缺失”这一细粒度问题。因此，本节将 \path{defect_outputs/visuals/}、\path{_analysis_detect/} 和 \path{defect_outputs/visuals/07_review_snapshots/} 中的关键资产作为完整前期实验写入报告。

\subsection{视频与运行配置}
根据 \path{dino/outputs/video_a/metadata/video_meta.json}，\path{20241031_222226.mp4} 的分辨率为 $1920\times1080$，帧率约 $29.851$ fps，共 $14108$ 帧，时长约 $472.609$s。早期粗 ROI 检测配置来自 \path{defect_outputs/model/run_config.json}，其核心参数见表~\ref{tab:video-a-config}。

\begin{table}[H]
\centering
\caption{\code{20241031\_222226.mp4} 粗 ROI / PatchCore 风格实验配置}
\label{tab:video-a-config}
\small
\begin{tabularx}{\textwidth}{L{0.34\textwidth}Y}
\toprule
项目 & 设置 \\
\midrule
输入视频 & \path{20241031_222226.mp4} \\
视频规格 & $1920\times1080$，约 $29.851$ fps，约 $472.609$s \\
检测 ROI & $x=600$--$1180$，$y=280$--$720$ \\
虚拟门线 & \code{gate\_x=900} \\
目标半径范围 & \code{min\_radius=65}，\code{max\_radius=135} \\
Patch crop 尺寸 & $224\times224$ \\
PatchCore 网格 & \code{grid=14}，记忆库上限 \code{max\_memory\_patches=4096} \\
训练正常段 & $0$--$350$s，步长 $1.0$s \\
验证正常段 & $350$--$360$s \\
扫描区间 & $360$s 到视频结束，步长 $0.5$s \\
事件规则 & 至少 $2$ 个观测，事件间隔阈值 $1.25$s \\
运行阈值 & \code{threshold=4.0}，来自 \code{minimum\_operational\_threshold} \\
\bottomrule
\end{tabularx}
\end{table}

\subsection{PatchCore 实现细节与改进}
\path{defect_detection_pipeline.py} 中的 \code{TexturePatchCore} 是一个面向本视频场景重写的轻量 PatchCore 风格实现。标准 PatchCore 通常使用 ImageNet 预训练 CNN 提取多层 patch embedding，再建立正常样本 patch memory bank，并用最近邻距离度量异常程度。本项目早期没有稳定的大规模标签，也希望保持 CPU 可运行和可解释性，因此没有直接使用重型深度特征，而是将 PatchCore 思想改造为“手工 patch 特征 + 正常记忆库 + 最近邻异常分数”的工业原型。

具体实现过程如下：
\begin{enumerate}[label=\arabic*.]
    \item 目标定位：先在手工 ROI $[600,280,1180,720]$ 内检测圆形目标，并在 \code{gate\_x=900} 附近形成虚拟检测门线，得到中心 $(c_x,c_y)$、半径 $r$ 和检测置信度。
    \item Crop 构造：以检测到的圆形目标为中心，按半径 $1.22$ 倍扩展，裁剪并 resize 到 $224\times224$。
    \item Patch 网格：将 crop 划分为 $14\times14$ 网格。为避免背景边缘 patch 干扰，只保留距离 crop 中心不超过 \code{inner\_radius\_px=72} 的内部 patch，即 \code{patch\_mask}。
    \item Patch 特征：每个有效 patch 提取 $18$ 维手工特征，包括 HSV hue 的 $\sin/\cos$、饱和度和亮度均值/方差、Lab 均值/方差、Sobel 梯度、Laplacian、灰度标准差、黄色区域比例以及分位数统计。
    \item 正常记忆库：从 $0$--$350$s 正常训练段收集 $343$ 个 normal crops，提取所有有效 patch 特征后用 \code{StandardScaler} 标准化。若 patch 数超过 \code{max\_memory\_patches=4096}，则用 \code{MiniBatchKMeans} 压缩为 $4096$ 个记忆中心；否则直接保存标准化 patch 特征。
    \item 最近邻评分：推理时对测试 crop 提取同样 patch 特征，使用 \code{NearestNeighbors} 计算每个 patch 到正常记忆库的欧氏最近邻距离。图像级 PatchCore score 取距离最高的前 $8\%$ patch 均值，强调局部异常而不是全图平均异常。
    \item 热力图输出：将 patch 距离写回 $14\times14$ 网格，生成异常热力图，用于 \path{04_defect_gallery.jpg}、review 视频和高级分析图。
\end{enumerate}

相较于朴素 PatchCore，本项目做了三点面向工业场景的改进。第一，使用内部圆形 patch mask，只保留目标端面附近 patch，减少大 ROI 中背景和边缘噪声。第二，使用黄色比例、饱和度和纹理强度构造 hybrid score：\code{hybrid\_score = patchcore\_score + color\_texture\_penalty}。当目标区域黄色/米黄色证据不足、饱和度和纹理异常时，分数会被保守抬高。第三，使用事件级过滤，只有连续超过阈值的观测才被归入持续异常事件，孤立高分样本不直接作为最终缺陷事件。

这些改进的动机是明确的：单纯最近邻距离容易受目标定位抖动、光照变化和局部反光影响；而工业缺陷检测更关心持续、可复核的异常证据。因此早期算法将 PatchCore 的局部异常思想与颜色/纹理先验、虚拟门线定位和事件过滤结合，形成可解释的视频异常检测原型。但它仍然属于粗 ROI 异常检测，后续才进一步收敛到 side micro ROI 缺失检测。

\subsection{初步定位与诊断观察帧}
\path{_analysis_detect/} 保存了该视频多个时间点的检测观察图，覆盖正常起始段、中段、事件开始附近和事件内部。图~\ref{fig:video-a-analysis-detect} 将 $5$s、$60$s、$300$s、$355$s、$374$s、$390$s 和 $420$s 的检测图放在一起，用于展示粗 ROI 检测在时间轴上的状态变化。$374$s 以后正好进入 \path{defect_outputs/defect_events.csv} 记录的持续异常事件，因此这些图是解释早期事件定位的关键素材。

\begin{figure}[H]
\centering
\includegraphics[width=0.24\textwidth]{_analysis_detect/frame_0005s.jpg}
\includegraphics[width=0.24\textwidth]{_analysis_detect/frame_0060s.jpg}
\includegraphics[width=0.24\textwidth]{_analysis_detect/frame_0300s.jpg}
\includegraphics[width=0.24\textwidth]{_analysis_detect/frame_0355s.jpg}\\[3pt]
\includegraphics[width=0.24\textwidth]{_analysis_detect/frame_0374s.jpg}
\includegraphics[width=0.24\textwidth]{_analysis_detect/frame_0390s.jpg}
\includegraphics[width=0.24\textwidth]{_analysis_detect/frame_0420s.jpg}
\caption{\code{20241031\_222226.mp4} 早期检测观察帧。该组图来自 \code{\_analysis\_detect/}，用于复核粗 ROI / 虚拟门线检测在正常段、事件开始附近和事件内部的空间位置。}
\label{fig:video-a-analysis-detect}
\end{figure}

\subsection{粗 ROI 模型输出与核心可视化}
该实验的主输出位于 \path{defect_outputs/}。摘要文件显示，训练阶段收集了 $343$ 个 normal crops，验证阶段收集了 $11$ 个 validation crops；扫描阶段输出 $218$ 条观测，其中 $127$ 条被判为事件内缺陷观测，$91$ 条为非缺陷观测。最终事件从 $374.0$s 持续到 $439.5$s，持续 $65.5$s，峰值时间为 $397.0$s，峰值 hybrid score 为 $5.530$，峰值 PatchCore score 为 $3.888$。该事件的峰值饱和度为 $0.175$，峰值 yellow ratio 为 $0.123$，说明早期判别高度依赖颜色/纹理异常组合。

\begin{table}[H]
\centering
\caption{\code{20241031\_222226.mp4} 粗 ROI 实验输出统计}
\label{tab:video-a-results}
\small
\begin{tabularx}{\textwidth}{L{0.32\textwidth}L{0.18\textwidth}Y}
\toprule
Output Type & Count / Value & Notes \\
\midrule
Normal memory crops & 343 & 来自 $0$--$350$s 训练正常段。 \\
Validation crops & 11 & 来自 $350$--$360$s 验证正常段。 \\
Scanned observations & 218 & 扫描区间 $360$s 到视频结束，$0.5$s 步长。 \\
Non-defect observations & 91 & \path{observations.csv} 中 \code{is\_defect=False}。 \\
Event defect observations & 127 & \path{observations.csv} 中 \code{is\_defect=True}。 \\
Detected events & 1 & 来自 \path{defect_events.csv}。 \\
Event range & $374.0$--$439.5$s & 持续 $65.5$s。 \\
Peak hybrid score & 5.530 & 阈值为 $4.0$。 \\
Peak PatchCore score & 3.888 & 峰值时间 $397.0$s。 \\
\bottomrule
\end{tabularx}
\end{table}

图~\ref{fig:video-a-core-visuals} 展示该实验的核心可视化资产：粗 ROI 标定图、normal memory bank 样本、异常分数时间线、缺陷 gallery 和特征空间散点。它们构成了早期实验报告的基本证据链。

\begin{figure}[H]
\centering
\includegraphics[width=0.49\textwidth]{defect_outputs/visuals/01_roi_calibration.jpg}
\includegraphics[width=0.36\textwidth]{defect_outputs/visuals/02_normal_memory_mosaic.jpg}\\[4pt]
\includegraphics[width=0.72\textwidth]{defect_outputs/visuals/03_score_timeline.png}\\[4pt]
\includegraphics[width=0.42\textwidth]{defect_outputs/visuals/04_defect_gallery.jpg}
\includegraphics[width=0.42\textwidth]{defect_outputs/visuals/05_feature_space.png}
\caption{\code{20241031\_222226.mp4} 粗 ROI / PatchCore 风格实验核心可视化。从上到下依次包含 ROI 标定、正常记忆库样本、异常分数时间线、缺陷候选 gallery 和特征空间散点。}
\label{fig:video-a-core-visuals}
\end{figure}

\subsection{回放视频与关键帧复核}
\path{defect_outputs/visuals/06_annotated_diagnostic_preview.mp4} 是 $360$--$430$s 预览区间的带诊断面板视频，\path{defect_outputs/visuals/07_sampled_2fps_review.mp4} 是抽样 $2$ fps 检查视频。后者导出了 \path{07_review_snapshots/} 中的关键帧截图，包括事件前正常帧、事件开始帧、事件内部峰值附近帧和视频末尾正常帧。图~\ref{fig:video-a-review-snapshots} 将这些截图集中展示，说明早期检测输出并非只有 CSV 统计，也包含可复核的视频证据。

\begin{figure}[H]
\centering
\includegraphics[width=0.32\textwidth]{defect_outputs/visuals/07_review_snapshots/review_0360.00s_normal.jpg}
\includegraphics[width=0.32\textwidth]{defect_outputs/visuals/07_review_snapshots/review_0374.00s_defect.jpg}
\includegraphics[width=0.32\textwidth]{defect_outputs/visuals/07_review_snapshots/review_0382.50s_defect.jpg}\\[3pt]
\includegraphics[width=0.32\textwidth]{defect_outputs/visuals/07_review_snapshots/review_0397.00s_defect.jpg}
\includegraphics[width=0.32\textwidth]{defect_outputs/visuals/07_review_snapshots/review_0472.50s_normal.jpg}
\caption{\code{20241031\_222226.mp4} sampled review 关键帧。$360$s 和 $472.5$s 为正常参考，$374$s、$382.5$s 和 $397$s 位于持续异常事件内，用于检查事件是否具有视觉连续性。}
\label{fig:video-a-review-snapshots}
\end{figure}

\subsection{高级分析图与结论边界}
\path{defect_outputs/advanced_visuals/} 和 \path{defect_outputs/analysis/} 进一步给出了阈值分离、时间多指标、热力图 contact sheet、时空门线稳定性、无监督特征漂移、事件取证板和累计证据监控等分析图。根据 \path{advanced_metrics.json}，accepted normal 的 $95$ 分位 hybrid score 为 $2.817$，event defect 的最小分数为 $4.134$，两者间隔为 $1.318$；样本 $t=382.5$s 的阈值 margin 为 $+0.418$。这些指标说明粗 ROI 实验在 \path{20241031_222226.mp4} 上形成了较清楚的事件分数分离。

但这一结论仍然有边界：该视频实验检测的是较大的圆形/门控 ROI 的整体异常，而不是最终任务中的侧边接口米白色套/盖缺失。它能够说明“异常时段可以被粗 ROI 异常检测捕获”，却不能证明“side micro ROI 缺失检测已经准确”。因此，本报告将该视频作为完整的前期实验和失败/转向证据纳入正文，而把最终方法建立在 \path{20241101_161258.mp4} 的 side micro ROI 校准样本和固定上下检测带上。
\newpage
\section{数据与 ROI 校准}
\subsection{原始视频与抽帧策略}
最终 side micro ROI 检测使用的视频为 \path{20241101_161258.mp4}。根据 \path{cap_annotation_tool/tasks.json}，该视频帧率为 $29.953$ fps，原始帧大小为 $1920\times1080$。根据 showcase 输出摘要，视频结尾时间约为 $1091.209$s。ROI 校准工具抽取 $22$ 帧右半边裁剪图，裁剪图尺寸为 $960\times1080$，但保存的对象坐标会加回 \code{offset\_x=960}。

抽帧时间覆盖三类样本：第一类是 $17$s 左右的正常参考帧；第二类是 $120$s、$184$s、$234.5$s、$280$s、$360$s、$577.5$s、$759$s、$900$s、$1009.5$s 等中段疑似或复核帧；第三类是 $1072$s--$1084$s 的后段稀缺异常密集帧。这样设计的原因是缺陷样本稀少，若只均匀采样，训练集中几乎无法得到足够 \code{MISSING} patch。

\subsection{ROI 校准工具设计}
ROI 校准工具的 README 明确说明：前端只显示原视频右半边裁剪图，保存坐标时自动加上 \code{offset\_x=960}。前端支持快捷键 \code{1/2/3} 标记帧状态为 \code{NORMAL/DEFECT/UNKNOWN}，支持手动画框、删除框、保存 JSON，并可导出 CSV。该设计避免操作者在全图中被左侧无关区域干扰，同时保留了与原视频坐标一致的输出。

\subsection{ROI 样本类别与坐标系统}
ROI 样本 schema 包括以下层次：
\begin{enumerate}[label=\arabic*.]
    \item 帧级标签：\code{NORMAL}、\code{DEFECT}、\code{UNKNOWN}；
    \item 部位级缺失标签：\code{top}、\code{bottom}、\code{side\_upper}、\code{side\_lower}；
    \item 框类型：\code{sample}、\code{fixed\_region}、\code{top}、\code{bottom}、\code{side\_upper}、\code{side\_lower}、\code{ignore}；
    \item 框状态：\code{UNLABELED}、\code{PRESENT}、\code{MISSING}、\code{UNKNOWN}、\code{IGNORE}。
\end{enumerate}
在最终 ROI 校准文件中，所有有效训练框均为 \code{type=sample}，其状态为 \code{PRESENT} 或 \code{MISSING}。虽然 schema 支持部位级缺失勾选，但当前 \path{annotations.json} 中 \code{missing\_parts} 为空，实际训练信号来自小框状态。

\subsection{ROI 样本统计}
表~\ref{tab:annotation-stats} 汇总 ROI 样本统计。需要强调的是，\code{side\_roi\_labeled\_crops.csv} 中的 \code{frame\_state} 是按每个 patch 展开后的帧状态计数，因此 \code{NORMAL}=61、\code{DEFECT}=54 表示 patch 来源分布，不等同于帧数量。

\begin{table}[H]
\centering
\caption{ROI 样本统计}
\label{tab:annotation-stats}
\small
\begin{tabularx}{\textwidth}{L{0.26\textwidth}L{0.16\textwidth}L{0.18\textwidth}Y}
\toprule
Label Type & Count & Percentage & Notes \\
\midrule
校准帧总数 & 22 & 100.00\% & 来自 \path{tasks.json}，均为右半边裁剪图。 \\
帧级 \code{NORMAL} & 12 & 54.55\% & 来自 \path{annotations.json}。 \\
帧级 \code{DEFECT} & 9 & 40.91\% & 覆盖中后段异常/疑似异常帧。 \\
帧级 \code{UNKNOWN} & 1 & 4.55\% & 视野或证据不足。 \\
\code{sample} patch 总数 & 115 & 100.00\% & 全部有效训练框均为 \code{type=sample}。 \\
\code{PRESENT} patch & 99 & 86.09\% & 米白色套/盖可见。 \\
\code{MISSING} patch & 16 & 13.91\% & 米白色实体证据缺失。 \\
\code{UNKNOWN} patch & 0 & 0.00\% & 当前有效训练 patch 未使用 UNKNOWN。 \\
上侧检测带 patch & 58 & 50.43\% & \path{side_roi_training_summary.json} 中 upper。 \\
下侧检测带 patch & 57 & 49.57\% & \path{side_roi_training_summary.json} 中 lower。 \\
\bottomrule
\end{tabularx}
\end{table}

ROI patch 的尺寸分布为：宽度最小 $54$、平均 $83.17$、最大 $142$；高度最小 $48$、平均 $72.41$、最大 $112$。上侧 ROI 框的 $x$ 范围为 $1036$--$1901$、$y$ 范围为 $351$--$491$；下侧 ROI 框的 $x$ 范围为 $1040$--$1912$、$y$ 范围为 $837$--$956$。这直接支持“固定上下两个水平检测带”的定义。

\subsection{ROI 校准可视化}
图~\ref{fig:annotation-overlay} 和图~\ref{fig:side-roi-crops} 是 ROI 校准阶段的重要证据。前者展示 ROI 框在原始视频右半区域中的位置，后者展示裁剪出的 side micro ROI patch。两者说明最终检测器关注的是很小的侧边接口，而非整根样本。

\begin{figure}[H]
\centering
\includegraphics[height=0.78\textheight,keepaspectratio]{cap_annotation_tool/annotation_overlay_mosaic.jpg}
\caption{ROI 校准框叠加 mosaic。该图由 ROI 样本分析阶段生成，展示不同时间帧中 ROI 小框在右半边裁剪图中的位置，证明有效样本集中在上下侧边接口窄带。}
\label{fig:annotation-overlay}
\end{figure}

\begin{figure}[H]
\centering
\includegraphics[width=0.80\textwidth]{cap_annotation_tool/side_roi_crops_mosaic.jpg}
\caption{side micro ROI crop mosaic。图中按 \code{PRESENT/MISSING} 展示 ROI 裁剪 patch，是训练小 patch 分类器和复核 ROI 语义的核心素材。}
\label{fig:side-roi-crops}
\end{figure}

\section{方法}
\subsection{总体流程}
最终算法流程如下：
\begin{enumerate}[label=\arabic*.]
    \item 读取原始视频与 ROI 校准样本；
    \item 从 \code{sample} 框中构造 \code{PRESENT/MISSING} 训练样本；
    \item 在上下固定检测带中随机采样背景 patch，构造 \code{BACKGROUND} 类；
    \item 对校准 patch 做翻转、亮度扰动、缺失样本锐化/模糊/平移增强；
    \item 提取 $127$ 维颜色、纹理、空间池化和边缘特征；
    \item 训练 \code{ExtraTreesClassifier} 三分类模型；
    \item 每 $0.5$s 扫描一帧，使用蓝色端盖锚点约束候选 $x$ 位置；
    \item 批量计算候选 patch 特征并批量 \code{predict\_proba}；
    \item 将模型概率与米白色实体证据规则融合；
    \item 对每个检测带做 NMS 与候选替换；
    \item 将持续缺失观测聚合为事件，孤立缺失标记为 \code{FILTERED\_SPIKE}；
    \item 输出 \path{observations.csv}、\path{events.csv}、\path{summary.json} 和可视化视频。
\end{enumerate}

\subsection{Micro-ROI 定义}
最终 micro ROI 的定义来自 ROI 样本统计，而不是拍脑袋设置。上侧检测带为 $y=350$--$495$，下侧检测带为 $y=835$--$960$。默认候选窗口为 $84\times68$，与 ROI 样本平均尺寸接近。检测仅在原始视频右半区域 $x\geq960$ 内进行，并设置 \code{min\_candidate\_x=1000} 避免贴近裁剪边界的半截候选。

使用 micro ROI 的原因有三点。首先，缺陷视觉证据只存在于接口小区域，整根样本包含过多无关信息。其次，透明接口和金属反光在大 ROI 中会显著增加误判概率。第三，小 ROI 使 ROI 校准成本低，能够在缺陷样本稀缺时快速构建可用训练集。

\subsection{候选生成}
候选生成复用 \path{cap_defect_detection_pipeline.py} 中的蓝色端盖锚点检测与配对函数。对每组上下端盖锚点，算法计算中心 $x$，并生成三个相对偏移候选：\code{candidate\_dx\_left=-92}、\code{candidate\_dx\_near\_left=-64}、\code{candidate\_dx\_right=192}。每个候选在上、下两个固定检测带各取一个窗口。若蓝色锚点检测失败，则在右半区域以较大步长生成稀疏候选作为后备。

这种策略显著区别于全宽密集滑窗。全宽滑窗不仅计算量大，而且会产生大量背景候选；而蓝盖约束候选直接将搜索限制在可能出现侧边接口的相对位置上，使模型在实际推理中只需要检查少量 patch。

\subsection{Patch 特征提取}
函数 \code{crop\_features} 将 patch resize 到检测窗口大小后提取 $127$ 维特征，具体包括：
\begin{enumerate}[label=\arabic*.]
    \item HSV 与 Lab 六个通道的均值、标准差、$10$ 分位数、$90$ 分位数，共 $24$ 维；
    \item 低饱和高亮米白色 mask、暗区域 mask、蓝色 mask 的整体比例，共 $3$ 维；
    \item $4\times4$ 网格空间池化，每格提取灰度均值、灰度标准差、饱和度均值、亮度均值和白色比例，共 $80$ 维；
    \item Sobel 梯度均值与标准差、Laplacian 方差、Canny 边缘密度，共 $4$ 维；
    \item $16$ bin 灰度直方图，共 $16$ 维。
\end{enumerate}
这些特征足以描述小 patch 中的颜色实体、透明边缘、局部纹理和背景差异，同时保持模型轻量。

\subsection{Patch 分类器}
最终分类器为 \code{sklearn.ensemble.ExtraTreesClassifier}，参数包括 \code{n\_estimators=160}、\code{class\_weight=balanced}、\code{max\_features=sqrt}、\code{random\_seed=17}、\code{n\_jobs=-1}。训练数据来自 ROI 校准样本和自动背景采样：最终增强后训练标签计数为 \code{PRESENT}=396、\code{MISSING}=160、\code{BACKGROUND}=254。这里的数量大于原始 $99/16$，原因是对原始 ROI patch 进行了数据增强；\code{BACKGROUND} 来自上下检测带中与 ROI 框低重叠的随机采样。

内部 holdout 划分使用 $25\%$ stratified split，报告精度为 $96.55\%$，混淆矩阵如表~\ref{tab:holdout-confusion}。该结果只能说明增强 patch 内部分布上模型可分，不应等同于独立视频级缺陷检测精度。

\begin{table}[H]
\centering
\caption{内部 holdout 混淆矩阵（增强 patch 数据，非独立视频级评估）}
\label{tab:holdout-confusion}
\small
\begin{tabular}{lccc}
\toprule
真实类别 / 预测类别 & \code{BACKGROUND} & \code{MISSING} & \code{PRESENT} \\
\midrule
\code{BACKGROUND} & 61 & 0 & 3 \\
\code{MISSING} & 0 & 36 & 4 \\
\code{PRESENT} & 0 & 0 & 99 \\
\bottomrule
\end{tabular}
\end{table}

\subsection{规则证据融合}
仅依靠分类器仍不足以处理透明接口误判问题。最终代码新增 \code{sleeve\_evidence}，在 patch 内部区域计算低饱和高亮 mask、最大连通白色实体比例和 Canny 边缘密度。若 \code{white\_ratio>=0.40}、\code{sleeve\_blob>=0.42}、\code{edge\_density<=0.165} 且模型 \code{p\_missing<0.25}，则规则将候选强化为 \code{PRESENT}，解释为 \code{solid\_sleeve}。若边缘密度较高且米白色实体不足，则强化为 \code{MISSING}，解释为 \code{transparent\_or\_bare\_port}。

该规则体现了本任务的关键视觉先验：正常证据是米白色实体套/盖，而不是透明接口轮廓。透明接口越清楚，反而可能意味着套/盖缺失。

\subsection{帧级决策与重叠框筛选}
每个采样时刻会得到上、下两个检测带中的若干候选。算法按 band 分组，对候选按 score 排序，并使用 \code{nms\_iou=0.18} 进行重叠抑制。为避免真实缺失框被旁边的正常框压掉，代码加入候选替换逻辑：强 \code{MISSING} 可以替换弱 \code{PRESENT}；反过来，若旧缺失框 \code{p\_missing} 不高且新候选具有强 \code{solid\_sleeve} 证据，也允许恢复为 \code{PRESENT}。每个 band 最多保留 $4$ 个候选。

\subsection{事件级聚合}
函数 \code{apply\_event\_filter} 将 patch 级 \code{MISSING} 观测按近似 $x$ 位置和 band 分组。如果相邻缺失观测时间间隔不超过 \code{event\_gap\_sec=1.25}s，则归为同一事件。事件至少需要 \code{min\_event\_observations=2} 个观测；不足两次的孤立缺失会被标记为 \code{FILTERED\_SPIKE}。因此最终报警逻辑并不依赖单帧，而强调持续证据。

\subsection{可视化流水线}
\path{side_micro_roi_detector.py} 生成 \path{01_micro_score_timeline.png} 和 \path{02_micro_review_2fps.mp4}。其中检测视频会在原始帧上绘制上下固定检测带、候选框、状态标签和左侧诊断面板。

\path{make_side_micro_showcase_video.py} 重新训练同一模型，并对最后 $60$s 以 $8$ fps 扫描和渲染。左侧面板包含实时风险仪表、\code{CLEAR/ALERT/WATCH} 状态、最大缺失分数、上下通道风险条、风险 sparkline、最近 $20$s band risk strip、top evidence crops 以及 crop heatmap。右侧保留原视频检测区域，绘制 \code{PRESENT/MISSING/FILTERED\_SPIKE} 颜色框。

\section{实验设置}
\begin{table}[H]
\centering
\caption{最终 side micro ROI 检测实验设置}
\label{tab:experimental-setup}
\small
\begin{tabularx}{\textwidth}{L{0.34\textwidth}Y}
\toprule
项目 & 设置 \\
\midrule
核心视频 & \path{20241101_161258.mp4} \\
视频尺寸与帧率 & $1920\times1080$，$29.953$ fps（来自 \path{tasks.json}） \\
视频时长 & 约 $1091.209$s（来自 showcase summary 的 \code{end\_time\_sec}） \\
校准帧数 & $22$ 帧右半边裁剪图 \\
ROI 校准 patch & $115$ 个：\code{PRESENT}=99，\code{MISSING}=16 \\
训练样本（增强后） & \code{PRESENT}=396，\code{MISSING}=160，\code{BACKGROUND}=254 \\
检测带 & upper: $y=350$--$495$；lower: $y=835$--$960$ \\
窗口尺寸 & $84\times68$ \\
特征维度 & $127$ 维手工颜色/纹理/空间池化/边缘/直方图特征 \\
分类器 & ExtraTrees，$160$ trees，balanced class weight，\code{n\_jobs=-1} \\
扫描间隔 & $0.5$s，即约 $2$ fps 检测 \\
候选约束 & 蓝色端盖配对锚点 + 三个固定 $x$ 偏移；失败时稀疏右半区候选 \\
概率阈值 & \code{present\_threshold=0.48}，\code{missing\_threshold=0.30}，\code{background\_threshold=0.56} \\
事件规则 & 同位置同 band，间隔 $\leq1.25$s，至少 $2$ 次观测 \\
全视频输出 & $12{,}762$ 条观测、$92$ 个事件 \\
检查视频 & \path{side_micro_outputs/visuals/02_micro_review_2fps.mp4}，$2$ fps \\
Showcase 视频 & 最后 $60$s，$1031.209$s--$1091.209$s，$8$ fps，$481$ 帧 \\
\bottomrule
\end{tabularx}
\end{table}
\newpage
\section{结果}
\subsection{检测输出统计}
表~\ref{tab:detection-output-stats} 汇总最终全视频检测输出。全视频共 $2183$ 个采样时刻，每个时刻保留若干上/下通道候选，最终得到 $12{,}762$ 条观测。这里的 \code{MISSING} 是事件过滤后的持续缺失观测，\code{FILTERED\_SPIKE} 是未满足事件持续性要求的孤立缺失。

\begin{table}[H]
\centering
\caption{检测输出统计}
\label{tab:detection-output-stats}
\small
\begin{tabularx}{\textwidth}{L{0.28\textwidth}L{0.18\textwidth}Y}
\toprule
Output Type & Count & Notes \\
\midrule
全视频采样时刻 & 2183 & \path{observations.csv} 中唯一 \code{time\_sec} 数量，范围 $0.0$s--$1091.0$s。 \\
Patch 观测总数 & 12,762 & 每个采样时刻上、下 band 的候选输出。 \\
\code{PRESENT} & 12,204 & 事件过滤后正常实体证据。 \\
\code{MISSING} & 428 & 进入持续事件的缺失观测。 \\
\code{FILTERED\_SPIKE} & 130 & 原始缺失候选，但未形成至少两次持续观测。 \\
\code{UNKNOWN} & 0 & 最终检测输出中未保留 UNKNOWN。 \\
事件总数 & 92 & 来自 \path{side_micro_outputs/events.csv}。 \\
Showcase 观测 & 2,760 & 最后 $60$s、$8$ fps 输出。 \\
Showcase 事件 & 30 & 来自 \path{03_last60s_8fps_showcase_events.json}。 \\
\bottomrule
\end{tabularx}
\end{table}

\subsection{事件统计}
\begin{table}[H]
\centering
\caption{事件统计}
\label{tab:event-stats}
\small
\begin{tabularx}{\textwidth}{L{0.20\textwidth}L{0.18\textwidth}L{0.16\textwidth}L{0.22\textwidth}Y}
\toprule
Number of Events & Mean Duration & Max Risk & Affected Channel & Notes \\
\midrule
92（全视频） & $2.09$s & $0.8375$ & upper: 62；lower: 30 & \path{events.csv}，事件持续时间范围 $0.5$s--$9.5$s。 \\
30（最后 $60$s showcase） & $1.63$s & $0.8750$ & upper: 18；lower: 12 & \path{03_last60s_8fps_showcase_events.json}，用于演示与复核。 \\
\bottomrule
\end{tabularx}
\end{table}

\subsection{准确率与推理速度分析}
课程报告通常需要给出每张视频准确率、多视频平均准确率和推理速度。需要特别说明的是，当前项目文件夹中没有两段视频的完整逐帧 ground truth，也没有事件级真值文件；\path{dino/outputs/video_a/metrics/baseline_metrics.json} 也明确记录 \code{status=missing\_labels}。因此，本报告不能把模型自身输出、文件名中的 review 状态或抽查截图当作真实标签来计算严格 accuracy。表~\ref{tab:accuracy-speed} 将“严格可计算指标”和“当前可报告指标”分开列出。

\begin{table}[H]
\centering
\caption{每个视频的准确率可计算性与推理速度}
\label{tab:accuracy-speed}
\small
\begin{tabularx}{\textwidth}{L{0.22\textwidth}L{0.18\textwidth}L{0.20\textwidth}L{0.18\textwidth}Y}
\toprule
Video & Strict Accuracy & Available Metric & Inference Speed & Notes \\
\midrule
\path{20241031_222226.mp4} & TODO：缺少逐帧/事件级标签 & 粗 ROI 事件分离：accepted normal p95 $2.817$，event defect min $4.134$；检测到 $1$ 个持续事件 & PatchCore 粗 ROI 未保存逐帧 latency；DINO/颜色 baseline 为 $16.29$ ms/frame，约 $61.40$ FPS & 不能将检测器输出事件反过来当 ground truth。需要补充事件标签后计算 accuracy。 \\
\path{20241101_161258.mp4} & TODO：缺少全视频逐帧/事件级标签 & Patch 分类器内部 holdout accuracy $96.55\%$；全视频输出 $92$ 个事件 & side micro scan 日志约 $7.67$ FPS；review 渲染约 $11.43$ FPS；showcase 渲染约 $9.46$ FPS & $96.55\%$ 仅是增强 patch 内部划分结果，不等同于视频级准确率。 \\
多视频平均 & TODO & 不计算严格平均准确率 & 可报告速度均值需统一任务后再算 & 两个视频任务定义不同：一个是粗 ROI 异常检测，一个是 side micro ROI 缺失检测，不能直接平均 accuracy。 \\
\bottomrule
\end{tabularx}
\end{table}

如果后续补充真值标签，建议使用统一的事件级标签表，例如 \code{video,start\_sec,end\_sec,label}，将预测事件与真值事件按时间 IoU 匹配，计算 event-level precision、recall、F1 和 video-level accuracy。对于逐帧 accuracy，则需要对每个采样时刻给出 \code{OK/NG}，再与 \path{observations.csv} 或 \path{events.csv} 聚合后的帧级状态对齐。

\subsection{关键版本对比}
\begin{longtable}{L{0.18\textwidth}L{0.19\textwidth}L{0.20\textwidth}L{0.20\textwidth}L{0.15\textwidth}}
\caption{关键版本对比}
\label{tab:version-comparison}\\
\toprule
Version & Detection Region & Candidate Strategy & Main Problem & Improvement \\
\midrule
\endfirsthead
\toprule
Version & Detection Region & Candidate Strategy & Main Problem & Improvement \\
\midrule
\endhead
Coarse ROI baseline & 整体目标或较大手工 ROI & 圆/门控区域或大 ROI 异常分数 & 背景、透明件和反光混入，不能定位侧边接口小缺陷 & 产生早期异常检测框架和时间线，但不适合作为最终判别区域。 \\
Dense sliding-window micro ROI & 侧边窄带 & 全宽滑窗候选 & 候选数量多、背景 patch 多、速度慢 & 证明必须进一步约束候选位置。 \\
Blue-cap constrained candidates & 上下固定水平带 & 蓝色端盖配对锚点 + 固定偏移 & 锚点失败时仍可能漏候选 & 显著减少全宽搜索，并使候选贴近真实接口。 \\
Batch prediction optimized detector & 上下固定水平带 & 批量特征矩阵 + 批量 \code{predict\_proba} & 逐 patch 推理开销过高 & 将候选特征堆叠后一次性预测，提高全视频扫描效率。 \\
Rule-enhanced final detector & side micro ROI & 分类概率 + 米白色实体证据规则 & 透明接口可能被当成正常，NMS 可能压掉缺失框 & 增加 \code{white\_ratio/sleeve\_blob/edge\_density} 规则和候选替换逻辑。 \\
Showcase visualization version & 最后 $60$s 视频 & 同最终检测器，$8$ fps 渲染 & 单纯 CSV 不利于展示和复核 & 生成风险仪表、曲线、证据 crop 和热力图面板。 \\
\bottomrule
\end{longtable}

\subsection{失败实验与修复}
\begin{longtable}{L{0.18\textwidth}L{0.20\textwidth}L{0.20\textwidth}L{0.20\textwidth}L{0.14\textwidth}}
\caption{失败实验与修复}
\label{tab:failure-fix}\\
\toprule
Problem & Evidence & Cause & Fix & Remaining Risk \\
\midrule
\endfirsthead
\toprule
Problem & Evidence & Cause & Fix & Remaining Risk \\
\midrule
\endhead
大 ROI 误检 & \path{defect_outputs/}、\path{cap_defect_outputs/} 的调试图和时间线 & 缺陷区域太小，大 ROI 混入背景与反光 & 改为上下两个固定 side micro ROI 检测带 & 若目标发生大幅位移，固定带可能需重新标定。 \\
DINO/颜色基线无正式结论 & \path{dino/outputs/video_a/metrics/baseline_metrics.json} 显示缺少 labels & 没有真值标签，且颜色比例不能表达侧边套/盖缺失 & 保留为旧实验，最终采用小 patch 监督 & 后续可在有真值后做跨方法对比。 \\
密集滑窗过慢 & smoke/speed 目录与最终候选约束修改 & 全宽搜索产生大量背景候选 & 使用蓝色端盖锚点和固定偏移生成候选 & 蓝盖定位失败时仍需后备策略。 \\
逐 patch \code{predict\_proba} 慢 & \path{side_micro_roi_detector.py} 中改为 \code{model.predict\_proba(np.vstack(...))} & Python 循环和 sklearn 单样本调用开销大 & 批量预测所有候选 patch & 特征提取仍是主要计算开销之一。 \\
\code{MISSING} 过于保守 & \path{review_micro_t1074p0.jpg} 等检查图 & 缺陷样本少，分类器倾向 \code{PRESENT} & 数据增强 + balanced class weight + 降低 missing 阈值 & 可能增加孤立误报，需要事件过滤。 \\
透明接口误判为 \code{PRESENT} & \path{debug_t1054p5_after_rule.jpg} 等规则调试图 & 模型把透明接口轮廓当作正常结构 & 新增米白色实体规则，透明/裸露接口触发缺失 & 极端反光仍可能模拟米白实体。 \\
重叠框压掉真实缺失 & NMS 修复后的 \path{debug_t1054p5_after_fix.jpg} 等 & 高分正常框与真实缺失框重叠 & 增加强缺失替换弱正常、强实体恢复弱缺失逻辑 & 多目标近距离重叠时仍需复核。 \\
\bottomrule
\end{longtable}

\section{可视化与展示}
\subsection{ROI 校准可视化}
图~\ref{fig:annotation-overlay} 与图~\ref{fig:side-roi-crops} 属于 ROI 校准阶段可视化。它们对应项目从大 ROI 转向小 ROI 的关键证据：ROI 框集中在上下侧边窄带，而不是包围整根样本。这些图进入正文，因为它们解释了最终方法的 ROI 定义来源。

\subsection{检测时间线}
图~\ref{fig:micro-timeline} 由 \path{side_micro_roi_detector.py} 输出，展示全视频采样时间上的最大缺失概率和事件区间。该图是定量输出的概览图，但仍需结合视频检查确认事件是否为真实缺陷。

\begin{figure}[H]
\centering
\includegraphics[width=\textwidth]{side_micro_outputs/visuals/01_micro_score_timeline.png}
\caption{最终固定 side micro ROI 检测时间线。黄色曲线表示采样时刻最大缺失分数，红色框表示事件聚合得到的持续缺失区间。}
\label{fig:micro-timeline}
\end{figure}

\subsection{2 fps 检查视频}
\path{side_micro_outputs/visuals/02_micro_review_2fps.mp4} 是全视频低帧率检查视频，输出帧率为 $2$ fps。视频在右半区域绘制上下固定检测带和 patch 框，左侧面板显示检测器名称、当前时间、固定 band 坐标、窗口尺寸、present/missing 数量以及“只把米白色实体套/盖作为正常证据”的判别说明。它用于快速浏览正常段与异常段，属于定性验证和调试资产。

\subsection{最后一分钟 8 fps 高级展示视频}
\path{side_micro_outputs/visuals/03_last60s_8fps_showcase.mp4} 是本项目的重要展示成果。根据 \path{03_last60s_8fps_showcase_summary.json}，视频范围为 $1031.209$s--$1091.209$s，$8$ fps，共 $481$ 帧，输出事件 $30$ 个。图~\ref{fig:showcase-snaps} 展示了三个关键帧截图。

\begin{figure}[H]
\centering
\includegraphics[width=0.32\textwidth]{side_micro_outputs/visuals/03_showcase_snapshot_05s.jpg}
\includegraphics[width=0.32\textwidth]{side_micro_outputs/visuals/03_showcase_snapshot_30s.jpg}
\includegraphics[width=0.32\textwidth]{side_micro_outputs/visuals/03_showcase_snapshot_55s.jpg}
\caption{最后 $60$s、$8$ fps showcase 视频关键帧。左侧面板展示实时风险、CLEAR/ALERT/WATCH 状态、风险条、曲线、证据 crop 和热力图；右侧保留原视频检测框。}
\label{fig:showcase-snaps}
\end{figure}

该视频的价值在于把 patch 级概率、规则证据、事件状态和原始视觉上下文同时展示出来。对于工业检测项目，可视化不仅是演示材料，也是发现误检、漏检和规则不合理的重要调试工具。

\section{误差分析}
\subsection{自动大 ROI 失败原因}
自动大 ROI 失败的根本原因是检测问题被错误地表述为“看整根样本是否异常”。在该表述下，算法确实可以得到纹理异常分数、颜色比例和目标定位框，但这些信号与“侧边接口米白色套/盖是否缺失”并不一一对应。正常样本中的透明件和金属反光可能产生高异常分数；缺陷样本中的真实证据则可能只占大 ROI 的很小一部分，被背景和主体结构淹没。

\subsection{透明接口与反光干扰}
透明接口是本任务最容易造成语义混淆的结构。若算法只学习边缘、形状或局部高亮，透明接口会被误认为有部件存在。实际缺陷定义要求观察米白色实体套/盖，因此透明接口的清晰可见反而可能是缺失证据。金属反光则会使低饱和高亮区域增加，可能伪造米白实体特征。因此最终规则同时检查白色比例、连通实体和边缘密度，而不是只看亮度。

\subsection{缺陷样本稀缺}
ROI 校准集中原始 \code{MISSING} patch 只有 $16$ 个。虽然数据增强扩展到 $160$ 个训练样本，但这些样本仍来自少量原始实例，不能代表所有可能的缺陷形态。内部 holdout 报告不能替代独立测试集；当前结果应视为工程可运行原型，而不是经过充分统计验证的生产模型。

\subsection{候选生成与计算效率}
单纯滑窗会产生大量背景 patch，并使 \code{predict\_proba} 调用次数迅速增加。最终方案通过蓝色端盖锚点和固定偏移减少候选数量，并使用批量 \code{predict\_proba} 提高效率。但该策略依赖蓝色端盖检测，一旦端盖被遮挡、曝光异常或目标型号变化，候选位置可能偏移。

\subsection{事件级过滤的必要性}
工业视频中单帧误报不可避免，尤其是在反光、遮挡、运动模糊或候选位置抖动时。若单帧直接报警，误报会显著增加。最终事件规则要求同位置同 band 的缺失观测至少出现 $2$ 次，且相邻间隔不超过 $1.25$s。该策略能过滤孤立 spike，但也可能丢弃持续时间极短的真实缺陷。

\subsection{当前可能失败的场景}
当前方法仍可能在以下场景失败：
\begin{enumerate}[label=\arabic*.]
    \item 强反光在 patch 内形成大块低饱和高亮区域，模拟米白套/盖；
    \item 目标侧边接口被遮挡或裁剪，导致候选 patch 证据不足；
    \item 样本位置大幅偏移，固定上下检测带失效；
    \item 蓝色端盖定位失败，使候选 $x$ 位置偏离真实接口；
    \item 缺陷区域比训练样本更小或形态不同；
    \item 相机角度、光照、曝光和新型号样本导致颜色阈值和手工特征失配；
    \item 没有独立事件级真值时，无法可靠估计 precision、recall 和漏检率。
\end{enumerate}

\section{讨论}
\subsection{为什么固定上下两个水平带更适合本任务}
本项目最终方案有效的根本原因是 ROI 定义与缺陷语义对齐。侧边米白色套/盖缺失不是全局外观异常，而是局部装配证据缺失。固定上下水平带将模型视野限制在最可能发生缺陷的区域，减少了无关结构带来的统计噪声。对于此类工业视觉问题，正确的 ROI 定义往往比更复杂的分类器更重要。

\subsection{ROI 小样本校准集的作用}
ROI 小样本校准集不仅提供训练样本，还纠正了项目对“sample 框”的理解。若没有这些小框，算法会继续试图自动定位大目标区域，从而持续受背景和反光影响。$22$ 帧、$115$ 个 patch 的样本量并不大，但它明确告诉模型和工程人员：应该看哪里、哪些证据代表正常、哪些证据代表缺失。

\subsection{视觉先验与规则增强的价值}
本项目没有盲目扩大模型规模，而是把任务先验转化为规则：正常需要米白色实体，透明接口不是正常证据。这种先验与 ExtraTrees 分类器结合后，能够在小样本条件下获得更可解释的输出。规则增强也使调试更直接：当误检发生时，可以查看 \code{white\_ratio}、\code{sleeve\_blob} 和 \code{edge\_density} 是否符合视觉常识。

\subsection{工程可部署性}
最终方案使用 OpenCV、NumPy、scikit-learn 和轻量手工特征，不依赖 GPU。候选数量少，模型为树集成分类器，具备较好的部署可行性。输出 CSV/JSON 和视频可视化也便于与复核流程结合。若要进入生产部署，还需要增加跨视频验证、在线相机标定、报警去抖、复核闭环和生产日志记录。

\section{思考与改进}
\subsection{最难检测的情况}
本项目最难检测的情况不是视觉上完全消失的明显缺陷，而是“看起来像有结构、但缺少米白色实体证据”的透明接口。透明件会产生边缘、折射和局部高亮，粗 ROI PatchCore 会把它视为局部纹理异常，side micro ROI 分类器也可能把它误认为正常接口结构。另一个困难情况是强反光：反光区域可能形成低饱和高亮块，短时间内看起来很像米白色套/盖，从而压低缺失风险。

从两个视频的实验可以看出，难点还有尺度差异。对于 \path{20241031_222226.mp4}，粗 ROI 异常检测容易得到连续高分事件，但它并不能精确解释侧边小接口的装配状态；对于 \path{20241101_161258.mp4}，最终缺陷证据集中在 $80\times70$ 左右的小 patch 中，算法必须同时解决候选位置、局部分类、透明接口规则和事件过滤。也就是说，最难的问题并非“分类器不够强”，而是“目标语义和检测区域必须足够精确”。

\subsection{改进思路}
后续改进应按优先级推进：
\begin{enumerate}[label=\arabic*.]
    \item 建立事件级 ground truth。没有逐帧/事件标签时，无法给出严格准确率，也无法判断 $92$ 个 side micro 事件中哪些是真缺陷、哪些是误报。
    \item 将 \path{20241031_222226.mp4} 的粗 ROI 事件整理为 \code{OK/NG} 时间段，并将 \path{20241101_161258.mp4} 的 side micro 事件整理为 \code{PRESENT/MISSING} 事件，统一计算 event-level precision、recall 和 F1。
    \item 扩大 \code{MISSING} 样本。当前原始缺失 patch 只有 $16$ 个，虽然增强后训练为 $160$ 个，但本质上仍来自少数视觉形态。
    \item 用轻量 CNN 或 MobileNetV3 patch classifier 替代手工特征，并保留 \code{white\_ratio/sleeve\_blob/edge\_density} 作为可解释辅助特征。
    \item 引入目标跟踪或几何配准，使上下水平带能够随样本位置轻微漂移，而不是完全固定在当前相机坐标。
    \item 将透明接口作为独立 hard negative 类别加入训练，而不是只依赖规则把它推向 \code{MISSING}。
    \item 在线部署时记录每个事件的关键帧、top evidence crop、概率曲线和复核结果，形成主动学习闭环。
\end{enumerate}

\section{局限性}
当前报告必须明确以下局限：
\begin{enumerate}[label=\arabic*.]
    \item 没有独立测试集，内部 holdout 仅来自增强 patch；
    \item 没有逐帧或事件级 ground truth，无法计算严格 precision/recall；
    \item \code{MISSING} 原始 ROI 样本仅 $16$ 个，缺陷形态覆盖不足；
    \item 最终规则包含手工阈值，跨光照、跨相机、跨型号泛化能力未知；
    \item 固定水平带依赖当前相机视角和目标位置，生产中若机械位置变化需重新标定；
    \item showcase 视频是定性展示，不应被解释为完整验证。
\end{enumerate}

\section{结论}
本项目完成了一条面向工业视频侧边接口米白色套/盖缺失检测的完整工程路线。项目从自动大 ROI 和通用异常检测尝试出发，经过失败分析和需求反馈，最终将问题重新定义为上下两个固定 side micro ROI 中的小 patch 分类与事件统计。ROI 校准样本集提供了关键语义监督，最终检测器结合 ExtraTrees patch 分类器、米白色实体证据规则、候选约束、批量预测、NMS 修复和事件级过滤，输出全视频观测表、事件表、时间线、$2$ fps 检查视频和最后 $60$s 的 $8$ fps 高级 showcase 视频。

当前结果证明该方案具备清晰的工程逻辑和可复核的可视化输出，但尚不能证明生产级精度。后续工作应优先补充独立事件级真值，计算事件级 precision/recall，扩展异常样本，进行跨视频验证，并评估轻量 CNN 或 MobileNet 替代手工特征的收益。

\clearpage
\appendix
\section{Project Artifact Index}
\subsection{代码文件索引}
\begin{longtable}{L{0.30\textwidth}L{0.28\textwidth}L{0.20\textwidth}L{0.14\textwidth}}
\caption{代码文件索引}
\label{tab:code-index}\\
\toprule
Path & Function & Related Stage & Used in Report Section \\
\midrule
\endfirsthead
\toprule
Path & Function & Related Stage & Used in Report Section \\
\midrule
\endhead
\path{side_micro_roi_detector.py} & 最终固定上下水平带 side micro ROI 检测器，训练 ExtraTrees，扫描视频，输出观测/事件/时间线/2fps 视频。 & Final detector & Method, Results \\
\path{make_side_micro_showcase_video.py} & 生成最后 $60$s、$8$ fps 高级 showcase 视频和关键帧截图。 & Visualization & Visualization \\
\path{prepare_cap_annotation_set.py} & 从视频抽取 $22$ 帧右半边裁剪图并生成 ROI 校准任务。 & ROI setup & Data and ROI Calibration \\
\path{cap_annotation_tool/server.py} & 本地 ROI 校准前端 API，读取任务并保存 \path{annotations.json}。 & ROI calibration tool & Data and ROI Calibration \\
\path{cap_annotation_tool/app.js} & ROI 校准前端交互逻辑，支持帧标签、画框、坐标 offset 和导出。 & ROI calibration tool & Data and ROI Calibration \\
\path{cap_defect_detection_pipeline.py} & 端盖/较大 ROI 检测尝试，包含蓝色端盖定位、patch memory 和早期 side ROI 规则。 & Coarse/cap ROI & Pipeline Evolution \\
\path{defect_detection_pipeline.py} & 更早期粗 ROI / PatchCore 风格纹理异常检测管线。 & Coarse ROI baseline & Pipeline Evolution \\
\path{advanced_defect_analysis.py} & 针对早期粗 ROI 结果生成高级诊断图、事件分析和指标报告。 & Failure analysis & Error Analysis \\
\path{dino/dino_ad/*.py} & DINO/颜色基线框架，用于探索通用异常检测、评估与演示视频。 & Alternative baseline & Appendix \\
\path{patchcore-inspection-main/} & PatchCore 参考实现与模型资产。 & Reference baseline & Appendix \\
\bottomrule
\end{longtable}

\subsection{ROI 校准文件索引}
\begin{longtable}{L{0.30\textwidth}L{0.32\textwidth}L{0.16\textwidth}L{0.14\textwidth}}
\caption{ROI 校准文件索引}
\label{tab:annotation-index}\\
\toprule
Path & Description & Number of Samples & Notes \\
\midrule
\endfirsthead
\toprule
Path & Description & Number of Samples & Notes \\
\midrule
\endhead
\path{cap_annotation_tool/tasks.json} & ROI 校准任务文件，包含 $22$ 帧右半边裁剪图、时间、帧号、offset 和自动预测参考。 & 22 frames & \code{right\_half\_x=960} \\
\path{cap_annotation_tool/annotations.json} & ROI 校准主文件，包含帧级状态和 \code{sample} 小框。 & 22 frames / 115 boxes & 正文使用 \\
\path{cap_annotation_tool/side_roi_annotations.csv} & 展开后的 side ROI 校准表，含坐标和 band。 & 115 rows & upper/lower 统计 \\
\path{cap_annotation_tool/side_roi_labeled_crops.csv} & 有标签 crop manifest，指向 \path{side_roi_crops/}。 & 115 rows & 训练数据来源 \\
\path{cap_annotation_tool/side_roi_training_summary.json} & ROI 样本分析摘要，给出 \code{PRESENT/MISSING} 计数、band 范围和推荐固定检测带。 & 115 boxes & 关键依据 \\
\path{cap_annotation_tool/frames/} & ROI 校准用右半边裁剪帧。 & 22 images & 附录资产 \\
\path{cap_annotation_tool/side_roi_crops/} & ROI patch 裁剪图。 & 115 crops & 附录资产 \\
\bottomrule
\end{longtable}

\subsection{结果文件索引}
\begin{longtable}{L{0.30\textwidth}L{0.24\textwidth}L{0.22\textwidth}L{0.16\textwidth}}
\caption{结果文件索引}
\label{tab:result-index}\\
\toprule
Path & Output Type & Related Script & Key Metrics \\
\midrule
\endfirsthead
\toprule
Path & Output Type & Related Script & Key Metrics \\
\midrule
\endhead
\path{side_micro_outputs/summary.json} & 最终检测摘要 & \path{side_micro_roi_detector.py} & 12,762 observations；92 events \\
\path{side_micro_outputs/observations.csv} & patch 级观测表 & \path{side_micro_roi_detector.py} & \code{PRESENT}=12,204；\code{MISSING}=428；\code{FILTERED\_SPIKE}=130 \\
\path{side_micro_outputs/events.csv} & 事件表 & \path{side_micro_roi_detector.py} & 92 events；mean duration $2.09$s；max risk $0.8375$ \\
\path{side_micro_outputs/training_manifest.csv} & 原始训练 crop manifest & \path{side_micro_roi_detector.py} & 369 rows，未含增强副本 \\
\path{side_micro_outputs/side_micro_model.pkl} & 模型文件 & \path{side_micro_roi_detector.py} & ExtraTrees + config + train report \\
\path{side_micro_outputs/visuals/03_last60s_8fps_showcase_summary.json} & showcase 摘要 & \path{make_side_micro_showcase_video.py} & 481 frames；30 events \\
\path{side_micro_outputs/visuals/03_last60s_8fps_showcase_observations.csv} & showcase 观测表 & \path{make_side_micro_showcase_video.py} & 2,760 observations \\
\path{side_micro_outputs/visuals/03_last60s_8fps_showcase_events.json} & showcase 事件表 & \path{make_side_micro_showcase_video.py} & 30 events；max risk $0.875$ \\
\path{cap_defect_outputs/summary.json} & 早期端盖/大 ROI 结果摘要 & \path{cap_defect_detection_pipeline.py} & 5,576 observations；79 events \\
\path{defect_outputs/summary.json} & 更早期粗 ROI 结果摘要 & \path{defect_detection_pipeline.py} & 218 observations；1 event \\
\path{defect_outputs/model/run_config.json} & 粗 ROI 运行配置 & \path{defect_detection_pipeline.py} & ROI、训练/验证/扫描区间、阈值 \\
\path{defect_outputs/observations.csv} & 粗 ROI 逐帧观测 & \path{defect_detection_pipeline.py} & 218 rows；127 defect observations \\
\path{defect_outputs/defect_events.csv} & 粗 ROI 事件表 & \path{defect_detection_pipeline.py} & 374.0--439.5s；duration 65.5s \\
\path{defect_outputs/analysis/advanced_metrics.json} & 粗 ROI 高级分析 & \path{advanced_defect_analysis.py} & accepted normal p95 $2.817$；defect min $4.134$ \\
\path{dino/outputs/video_a/metrics/baseline_metrics.json} & DINO/颜色基线评估状态 & \path{dino/dino_ad/evaluate.py} & missing labels；无正式精度 \\
\bottomrule
\end{longtable}

\subsection{可视化文件索引}
\begin{longtable}{L{0.30\textwidth}L{0.16\textwidth}L{0.34\textwidth}L{0.12\textwidth}}
\caption{可视化文件索引}
\label{tab:visual-index}\\
\toprule
Path & Type & Description & Included Location \\
\midrule
\endfirsthead
\toprule
Path & Type & Description & Included Location \\
\midrule
\endhead
\path{cap_annotation_tool/annotation_overlay_mosaic.jpg} & JPG & ROI 校准框叠加图，展示小框集中于上下侧边接口。 & 正文 \\
\path{cap_annotation_tool/side_roi_crops_mosaic.jpg} & JPG & \code{PRESENT/MISSING} side ROI crop mosaic。 & 正文 \\
\path{side_micro_outputs/visuals/01_micro_score_timeline.png} & PNG & 最终全视频缺失分数时间线和事件区间。 & 正文 \\
\path{side_micro_outputs/visuals/03_showcase_snapshot_05s.jpg} & JPG & showcase 5s 关键帧。 & 正文 \\
\path{side_micro_outputs/visuals/03_showcase_snapshot_30s.jpg} & JPG & showcase 30s 关键帧。 & 正文 \\
\path{side_micro_outputs/visuals/03_showcase_snapshot_55s.jpg} & JPG & showcase 55s 关键帧。 & 正文 \\
\path{side_micro_outputs/visuals/debug_t1054p5_after_rule.jpg} & JPG & 规则增强调试图，展示透明/裸露接口处理。 & 附录 \\
\path{side_micro_outputs/visuals/debug_t1054p5_after_fix.jpg} & JPG & NMS/候选替换修复后调试图。 & 附录 \\
\path{cap_defect_outputs/visuals/01_blue_cap_locator.jpg} & JPG & 早期蓝色端盖定位可视化。 & 附录 \\
\path{cap_defect_outputs/visuals/03_score_timeline.png} & PNG & 早期端盖/大 ROI 分数时间线。 & 附录 \\
\path{_analysis_detect/frame_0005s.jpg--frame_0420s.jpg} & JPG & \code{20241031} 多时间点检测观察帧，覆盖正常段、事件开始和事件内部。 & 正文 \\
\path{defect_outputs/visuals/01_roi_calibration.jpg} & JPG & 粗 ROI 标定图。 & 正文 \\
\path{defect_outputs/visuals/02_normal_memory_mosaic.jpg} & JPG & 粗 ROI 正常记忆库样本 mosaic。 & 正文 \\
\path{defect_outputs/visuals/03_score_timeline.png} & PNG & \code{20241031} 异常分数时间线。 & 正文 \\
\path{defect_outputs/visuals/04_defect_gallery.jpg} & JPG & 粗 ROI 缺陷候选 gallery。 & 正文 \\
\path{defect_outputs/visuals/05_feature_space.png} & PNG & 粗 ROI 特征空间散点图。 & 正文 \\
\path{defect_outputs/visuals/07_review_snapshots/*.jpg} & JPG & sampled review 正常/缺陷关键帧截图。 & 正文 \\
\path{defect_outputs/advanced_visuals/08_threshold_separation_report.png} & PNG & 粗 ROI 阈值分离报告图。 & 附录 \\
\path{defect_outputs/advanced_visuals/09_multimetric_temporal_dashboard.png} & PNG & 粗 ROI 多指标时间仪表板。 & 附录 \\
\path{defect_outputs/advanced_visuals/10_patchcore_heatmap_contact_sheet.jpg} & JPG & PatchCore heatmap contact sheet。 & 附录 \\
\path{defect_outputs/advanced_visuals/11_spatiotemporal_gate_map.png} & PNG & 粗 ROI 时空门线稳定性图。 & 附录 \\
\path{defect_outputs/advanced_visuals/12_unsupervised_feature_drift_atlas.png} & PNG & 无监督特征漂移图谱。 & 附录 \\
\path{defect_outputs/advanced_visuals/13_event_forensics_board.jpg} & JPG & 粗 ROI 事件取证板。 & 附录 \\
\path{defect_outputs/advanced_visuals/14_cumulative_evidence_monitor.png} & PNG & 累计证据监控图。 & 附录 \\
\path{dino/outputs/video_a/figures/roi_contact_sheet.jpg} & JPG & DINO/颜色基线 ROI contact sheet。 & 附录 \\
\bottomrule
\end{longtable}

\subsection{视频文件索引}
\begin{longtable}{L{0.32\textwidth}L{0.12\textwidth}L{0.18\textwidth}L{0.20\textwidth}L{0.10\textwidth}}
\caption{视频文件索引}
\label{tab:video-index}\\
\toprule
Path & FPS & Time Range & Purpose & Notes \\
\midrule
\endfirsthead
\toprule
Path & FPS & Time Range & Purpose & Notes \\
\midrule
\endhead
\path{20241101_161258.mp4} & 29.953 & 约 0--1091.209s & 最终 side micro ROI 检测原始视频。 & 核心输入 \\
\path{20241031_222226.mp4} & 29.851 & 约 0--472.609s & 粗 ROI / PatchCore 风格前期完整实验视频。 & 正文前期实验 \\
\path{side_micro_outputs/visuals/02_micro_review_2fps.mp4} & 2 & 全视频 & 检查最终检测结果、框和诊断面板。 & 正文说明 \\
\path{side_micro_outputs/visuals/03_last60s_8fps_showcase.mp4} & 8 & 1031.209--1091.209s & 最后一分钟高级展示视频。 & 重点成果 \\
\path{cap_defect_outputs/visuals/06_review_2fps.mp4} & 2 & 全视频扫描输出 & 早期端盖/大 ROI 检查视频。 & 旧实验 \\
\path{defect_outputs/visuals/06_annotated_diagnostic_preview.mp4} & TODO & 360--430s 预览 & 粗 ROI PatchCore 诊断预览。 & 正文说明 \\
\path{defect_outputs/visuals/07_sampled_2fps_review.mp4} & 2 & 抽样检查片段 & 粗 ROI sampled review。 & 正文说明 \\
\path{dino/outputs/video_a/demo_15s.mp4} & 15 & 15s demo & DINO/颜色基线演示。 & 旧实验 \\
\path{dino/outputs/video_a/demo_advanced_last2min_8fps.mp4} & 8 & 最后 2 分钟 & DINO/颜色基线高级演示。 & 旧实验 \\
\bottomrule
\end{longtable}

\subsection{失败实验索引}
\begin{longtable}{L{0.22\textwidth}L{0.24\textwidth}L{0.24\textwidth}L{0.18\textwidth}L{0.06\textwidth}}
\caption{失败实验索引}
\label{tab:failure-index}\\
\toprule
Experiment & Problem & Evidence File & Fix & Status \\
\midrule
\endfirsthead
\toprule
Experiment & Problem & Evidence File & Fix & Status \\
\midrule
\endhead
Coarse ROI PatchCore & 大 ROI 与真实小缺陷不对齐。 & \path{defect_outputs/summary.json}、\path{visuals/01_roi_calibration.jpg} & 转向 side micro ROI 校准。 & Archived \\
Color/DINO baseline & 无真值标签，颜色比例不能直接表达套/盖缺失。 & \path{dino/outputs/video_a/metrics/baseline_metrics.json} & 不作为最终结果，只保留参考。 & Archived \\
Cap detector broad side ROI & 自动侧边 ROI 仍偏大，误把透明/背景结构纳入证据。 & \path{cap_defect_outputs/summary.json}、debug 图 & 改为固定上下微小检测带。 & Superseded \\
Dense micro sliding & 候选过多，背景干扰与速度问题。 & \path{side_micro_outputs_smoke*/} & 蓝盖锚点 + 固定偏移候选。 & Fixed \\
Transparent port false present & 透明接口被当作正常结构。 & \path{debug_t1054p5_after_rule.jpg} & 增加实体米白套/盖证据规则。 & Fixed \\
NMS suppresses missing & 高分 \code{PRESENT} 框压掉缺失框。 & \path{debug_t1054p5_after_fix.jpg} & 增加候选替换逻辑。 & Fixed \\
Single-frame alarms & 孤立反光造成 spike。 & \path{observations.csv} 中 \code{FILTERED\_SPIKE}=130 & 事件级过滤。 & Mitigated \\
\bottomrule
\end{longtable}

\section{Additional Visualizations}
除正文纳入的图像外，以下图像建议在正式汇报或论文版中补充展示：
\begin{enumerate}[label=\arabic*.]
    \item \path{side_micro_outputs/visuals/review_micro_t1074p0.jpg}、\path{review_micro_t1082p5.jpg}：检查关键异常时间点；
    \item \path{side_micro_outputs/visuals/debug_t1054p5_all_candidates.jpg}：候选生成与筛选调试；
    \item \path{cap_defect_outputs/visuals/debug_side_crop_mosaic.jpg}：早期 side ROI 尺度偏大的证据；
    \item \path{defect_outputs/advanced_visuals/08_threshold_separation_report.png} 至 \path{14_cumulative_evidence_monitor.png}：早期粗 ROI 分析图；
    \item \path{dino/outputs/video_a/figures/advanced_last2min_preview.jpg}：DINO/颜色基线演示截图。
\end{enumerate}

\section{Implementation Details}
\subsection{主要运行命令}
\begin{codeblock}
python defect_detection_pipeline.py run ^
  --video 20241031_222226.mp4 ^
  --out defect_outputs ^
  --train-end 350 ^
  --val-start 350 ^
  --val-end 360 ^
  --scan-start 360 ^
  --scan-step 0.5 ^
  --preview-start 360 ^
  --preview-end 430

python defect_detection_pipeline.py review-video ^
  --video 20241031_222226.mp4 ^
  --observations defect_outputs/observations.csv ^
  --model-dir defect_outputs/model ^
  --out defect_outputs/visuals/07_sampled_2fps_review.mp4 ^
  --fps 2 ^
  --sample-step 0.5 ^
  --snapshots-dir defect_outputs/visuals/07_review_snapshots

python prepare_cap_annotation_set.py

python side_micro_roi_detector.py run ^
  --video 20241101_161258.mp4 ^
  --annotations cap_annotation_tool ^
  --out side_micro_outputs ^
  --sample-step 0.5

python make_side_micro_showcase_video.py ^
  --video 20241101_161258.mp4 ^
  --annotations cap_annotation_tool ^
  --out side_micro_outputs/visuals/03_last60s_8fps_showcase.mp4 ^
  --seconds 60 ^
  --fps 8
\end{codeblock}

\subsection{主要依赖}
\path{requirements_defect_pipeline.txt} 列出核心依赖：\code{opencv-python-headless}、\code{numpy}、\code{pillow}、\code{scikit-learn}、\code{tqdm}。早期 DINO/AnomalyDINO 实验还涉及 PyTorch 与 DINOv2 本地权重，但最终 side micro ROI 检测器不依赖 GPU。

\subsection{复现实验的输入输出}
复现实验至少需要以下输入：\path{20241101_161258.mp4}、\path{cap_annotation_tool/tasks.json}、\path{cap_annotation_tool/annotations.json} 和 \path{cap_annotation_tool/frames/}。运行最终检测器后，应生成 \path{side_micro_outputs/summary.json}、\path{observations.csv}、\path{events.csv}、\path{side_micro_model.pkl}、\path{visuals/01_micro_score_timeline.png} 和 \path{visuals/02_micro_review_2fps.mp4}。运行 showcase 脚本后，应生成 \path{03_last60s_8fps_showcase.mp4}、对应 observations/events/summary 以及 $5$s、$30$s、$55$s 三张快照。

\section{Remaining TODOs}
\begin{table}[H]
\centering
\caption{剩余待确认事项}
\label{tab:todos}
\small
\begin{tabularx}{\textwidth}{L{0.34\textwidth}Y}
\toprule
TODO & Reason \\
\midrule
构建独立测试集 & 当前 holdout 来自增强 patch，不能替代跨时间/跨视频独立测试。 \\
补充事件级 ground truth & 需要明确每个事件是否真实缺陷，才能计算 event precision/recall。 \\
计算帧级与事件级 precision/recall & 当前仅能报告工程输出统计和定性/抽查验证。 \\
增加更多 \code{MISSING} 样本 & 原始缺失 patch 只有 $16$ 个，形态覆盖不足。 \\
进行跨视频验证 & 需要验证不同光照、相机角度、样本型号和产线状态下的泛化能力。 \\
加入稳定目标跟踪或在线标定 & 固定上下水平带对目标姿态变化敏感。 \\
尝试轻量 CNN / MobileNet patch 分类器 & 在数据扩充后比较手工特征与学习特征的收益。 \\
建立复核闭环 & 将 \code{FILTERED\_SPIKE}、高风险正常和低置信缺陷样本回流为后续训练样本。 \\
重新绘制正式论文图 & 当前部分调试图适合工程复核，论文版可进一步统一风格。 \\
\bottomrule
\end{tabularx}
\end{table}

\section*{致谢}
\addcontentsline{toc}{section}{致谢}
感谢金长龙老师在本项目选题、任务理解、实验推进与报告整理过程中给予的指导与帮助。老师对工业视觉检测场景、实验规范和结果表达的建议，使本项目能够从早期粗 ROI 尝试逐步收敛为更加清晰、可复核的微小 ROI 检测方案。谨此致谢。

\begin{thebibliography}{99}
\bibitem{opencv}
Bradski, G. The OpenCV Library. \emph{Dr. Dobb's Journal of Software Tools}, 2000.

\bibitem{sklearn}
Pedregosa, F., Varoquaux, G., Gramfort, A., et al. Scikit-learn: Machine Learning in Python. \emph{Journal of Machine Learning Research}, 12:2825--2830, 2011.

\bibitem{extratrees}
Geurts, P., Ernst, D., and Wehenkel, L. Extremely randomized trees. \emph{Machine Learning}, 63:3--42, 2006.

\bibitem{patchcore}
Roth, K., Pemula, L., Zepeda, J., et al. Towards Total Recall in Industrial Anomaly Detection. \emph{CVPR}, 2022.

\bibitem{mvtec}
Bergmann, P., Fauser, M., Sattlegger, D., and Steger, C. MVTec AD -- A Comprehensive Real-World Dataset for Unsupervised Anomaly Detection. \emph{CVPR}, 2019.
\end{thebibliography}

\end{document}
